% !PRINTED none
% [part-title / section-opening page: no running-head page number visible in any of the three bands; centred "Zweiter Abschnitt." / "Abel'sche Gruppen." title with rule, then \S.\,9 heading, then body text begins]
\begin{center}
\gesperrt{Zweiter Abschnitt.}

\gesperrt{Abel'sche Gruppen.}

\rule{0.3\linewidth}{0.4pt}

\S.\,9.

\gesperrt{Darstellung Abel'scher Gruppen durch eine Basis.}
\end{center}

Wir haben schon in \S.\,1 solche Gruppen, in denen bei der Composition das commutative Gesetz gilt, commutative oder \gesperrt{Abel'sche} Gruppen genannt. Bei diesen Gruppen, die in den Anwendungen von besonderer Wichtigkeit sind, herrschen viel einfachere Gesetze, als in den allgemeinen Gruppen, wie wir jetzt sehen werden. Es gelten für die Zusammensetzung der Elemente in diesen Gruppen dieselben Regeln, wie bei der Multiplication von Zahlen, und wir nennen demnach die Composita von Elementen einer solchen Gruppe, wie bei der Multiplication, Producte und Potenzen.

Auch hier beschränken wir uns auf die Betrachtung endlicher Gruppen.

Aus der Definition der Abel'schen Gruppen folgt, dass das commutative Gesetz auch bei der Composition der Theile einer solchen Gruppe gilt (\S.\,4). Jeder Theiler einer Abel'schen Gruppe ist selbst eine Abel'sche Gruppe, und ist zugleich Normaltheiler, so dass wir bei diesen Gruppen den Zusatz ,,Normal`` weglassen können.

Eine Gruppe, die nur aus den Wiederholungen eines Elementes besteht, wie
\[ 1, \; A, \; A^2 \; \ldots \; A^{a-1}, \]
haben wir schon früher (Bd. I, \S.\,148) eine cyklische Gruppe genannt, und wir behalten diese Bezeichnung hier bei. Wenn $a$ die kleinste positive Zahl ist, für die $A^a = 1$ ist, also $a$ der Grad der Gruppe, so heisst $a$ auch der Grad des Elementes $A$. Ist $m$ irgend ein Exponent, für den $A^m = 1$ ist, so ist $m$ ein
%CONT>>
